% !TeX root = ../../tesi.tex

For the integration of the Raspberry Pi Camera Module~3 Wide within the acquisition pipeline, the \textit{Picamera2} library was adopted.\cite{raspberrypi_picamera2_manual_2026,raspberrypi_picamera2_github_2026} Picamera2 is the Python interface built on top of the Raspberry Pi \textit{libcamera} stack and represents the modern replacement for the legacy Picamera interface.\cite{raspberrypi_picamera2_github_2026}

The library provides high-level access to camera initialization, stream configuration, frame acquisition, metadata handling, and camera controls such as exposure, gain, and autofocus.\cite{raspberrypi_picamera2_manual_2026} These features are particularly useful in embedded computer-vision applications, where the camera must be configured programmatically and integrated into a broader software pipeline.

In the present project, Picamera2 enabled direct access to RGB frames from Python, thereby simplifying the interaction between the camera and the processing software based on NumPy and OpenCV. This made it possible to implement image acquisition, parameter tuning, and frame extraction within the same software environment used for the rest of the pipeline.

Compared with lower-level camera handling, Picamera2 provides a practical abstraction that accelerates development while still exposing the controls needed for embedded deployment. This was particularly useful not only for the implementation of the acquisition pipeline itself, but also for the two-week automated data-collection phase, in which a stable and well-supported software interface was essential to manage continuous image acquisition and reliable data storage over extended periods. For this reason, Picamera2 represented a suitable software interface for Site~2, where ease of integration, reliability, and compatibility with the Raspberry Pi ecosystem were important design requirements.\cite{raspberrypi_picamera2_manual_2026,raspberrypi_picamera2_github_2026}
